%!TEX root=../../main.tex

In this section, we develop properties of CGL,
a generalized linear model which captures both the convex ReLU and convex gated
ReLU programs.
Let \( \calB = \cbr{b_1, \ldots, b_m} \) be a disjoint partition of the
feature indices \( [d] \).
Given regularization parameter \( \lambda \geq 0 \),
CGL solves the program:
\begin{equation}\label{eq:constrained-gl}
    \begin{aligned}
        p^*(\lambda) \!=\! \min_{w} \,
         & F_\lambda(w) \!:=\! \half \norm{X w \!-\! y}_2^2 + \lambda \sum_{\bi \in \calB} \norm{\wi}_2 \\
         & \quad \text{s.t.} \quad \Ki^\top \wi \leq 0  \text{ for all } \bi \in \calB,
    \end{aligned}
\end{equation}
where \( \Ki \in \R^{|\bi| \times a_\bi} \).
Solutions to \cref{eq:constrained-gl} are block sparse when \( \lambda \)
is sufficiently large, meaning \( \wi = 0 \) for a subset of blocks.
This is similar to the feature sparsity given by the lasso,
to which CGL naturally reduces when \( \bi = \cbr{i} \)
and \( \Ki = 0 \) for each \( \bi \in \calB \).
Although we consider squared-error, our results generalize to
strictly convex losses --- see \cref{app:general-losses} for comments.

The convex reformulations introduced in the previous section are instances
of CGL using the basis function \( X = [D_1 Z \ldots D_p Z] \),
where \( p = |\tilde \calD_Z| \).
For gated ReLU models, \( \Ki = 0 \) while ReLU models set
\( \Ki = - Z^\top (2 D_i - I) \).
For both problems, block sparsity from the group \( \ell_1 \) penalty
induces neuron sparsity in the final solution.

Our goal is to characterize the solution function of CGL,
\[
    \solfn(\lambda) := \argmin_{w \, : \, K_\bi^\top \wi \leq 0} F_\lambda(w)
\]
For a general data matrix, \( F_\lambda \) is not strictly convex and
CGL may admit multiple solutions --- these correspond to networks which
are not related by permutation.
As such, \( \solfn \) is a point-to-set map and
we must use a criterion to define a function;
for instance, the min-norm solution mapping
\[
    \wmin(\lambda) = \argmin \cbr{ \norm{w}_2 : w \in \solfn(\lambda) },
\]
defines a function for all \( \lambda \geq 0 \).

Now we introduce notation that will be used throughout this section.
Let \( \hat y(\lambda) = X w \) for \( w \in \solfn(\lambda) \) denote
the optimal model fit, which is the same for any choice of optimal \( w \)
(\cref{lemma:unique-fit-cgl}).
%Note that the sum of group norms is also constant over
%\( \solfn(\lambda) \) as a consequence.
Similarly, define the optimal residual \( r(\lambda) := y - \hat y(\lambda) \)
and \( \ci(\lambda) := \Xbi^\top r(\lambda) \) as the
correlation vector for block \( \bi \).
We write \( c \in \R^d \) for the concatenation
of these block-vectors.
%Since \( \hat y(\lambda) \) is the same for all optimal weights,
%\( r \) and \( c \) are constant over \( \solfn(\lambda) \).
Finally, let \( \ri \) be the dual parameters for the constraint
\( \Ki^\top \wi \leq 0 \), \( \rho \) their concatenation,
and \( K \) the block-diagonal matrix with blocks given by \( \Ki \).

The Lagrangian associated with \cref{eq:constrained-gl} is
\begin{equation}\label{eq:lagrangian-cgl}
    \calL(w, \rho) = \half \norm{X w - y}_2^2 + \lambda \sum_{\bi \in \calB} \norm{\wi}_2 + \abr{K \rho, w}.
\end{equation}
The constraints are linear, strong duality attains if feasibility holds,
and the necessary and sufficient conditions for primal-dual pair \( (w, \rho) \)
to be optimal are the KKT conditions:
\begin{equation}\label{eq:kkt-cgl}
    \begin{aligned}
        \Xbi^\top (X w - y) + \Ki \ri + s_\bi & = 0                                \\
        \Ki^\top \wi                          & \leq 0                             \\
        [\ri]_j \cdot [\Ki]_j^\top \wi        & = 0 \quad \forall \, j \in [a_\bi] \\
        \ri                                   & \geq 0,
    \end{aligned}
\end{equation}
where \( s_\bi \in \partial \lambda \norm{\wi}_2 \).
Since the KKT conditions hold for every combination of optimal primal-dual
pair \citep{boyd2014convex}, we always use the min-norm dual
optimal parameter \( \rmin \) with no loss of generality.
To simplify our notations, we define \( \vi := \ci - \Ki \rmin_\bi \).
In what follows, all proofs are deferred to \cref{app:cgl}.

\subsection{Describing the Optimal Set}

Stationary of the Lagrangian implies the equicorrelation set
\[
    \equi = \cbr{ \bi \in \calB : \norm{\vi}_2 = \lambda }.
\]
contains all blocks which may be active for fixed \( \lambda \).
That is, the active set \( \act(w) = \cbr{\bi : \wi \neq 0} \) satisfies
\( \act(w) \subseteq \equi \) for every \( w \in \solfn(\lambda) \).
However, not all blocks in \( \equi \) may be non-zero for some solution.
Thus, we define
\[
    \calS_\lambda = \cbr{\bi \in \calB : \exists w \in \solfn(\lambda), \wi \neq 0},
\]
which is the set of blocks supported by some solution.

Our first result combines the KKT conditions with uniqueness of
\( \hat y(\lambda) \) to characterize the solution set for fixed
\( \lambda > 0 \).
\begin{restatable}{proposition}{solFnCGL}\label{prop:sol-fn-cgl}
    Fix \( \lambda > 0 \).
    The optimal set for the CGL problem is
    given by
    \begin{equation}\label{eq:sol-fn-cgl}
        \begin{aligned}
            \solfn(\lambda) \!=\!
            \big\{ & w \! \in \! \R^d \!:\!
            \forall \, \bi \! \in \! \calS_\lambda,
            \wi \!= \! \alpha_\bi \vi, \alpha_\bi \!\geq\! 0, \,             \\
                   & \quad \forall \, b_j \in \calB \setminus \calS_\lambda,
            \w_{b_j} = 0, \, X w = \hat y
            \big\}
        \end{aligned}
    \end{equation}
\end{restatable}

Since this characterization is implicit due to the dependence
on \( \calS_\lambda \), we also give an alternative and explicit
construction in \cref{prop:cgl-alternate}, which shows that when \( K = 0 \)
we may replace \( \calS_\lambda \)
with \( \equi \);
We prefer \cref{prop:sol-fn-cgl} to \cref{prop:cgl-alternate} since it
better mirrors this simpler setting.
However, \cref{prop:cgl-alternate} can be substituted wherever desired.

Now that we know ``shape'' of the solution set, it is
possible to obtain simple conditions for existence of a unique solution.
As an immediate consequence of \cref{prop:sol-fn-cgl},
the solution map is a subset of directions in \( \Null(\Xe) \).
\begin{corollary}\label{cor:uniqueness-cond-cgl}
    If \( w, w' \in \solfn(\lambda) \) and \( z' = w - w' \), then
    \[
        z_\equi' \in \calN_\lambda \! := \! \Null(\Xe) \cap \cbr{z_\equi :  \forall \bi \in \equi, \, \zi \! = \! \alpha_\bi \vi}.
    \]
    As a result, the group lasso solution is unique if \( \calN_\lambda = \cbr{0} \).
\end{corollary}

\cref{cor:uniqueness-cond-cgl} extends a similar result for the lasso
to CGL \citep[Eq. 9]{tibshirani2013unique} and
implies the solution is unique for all \( \lambda \geq 0 \) if
the columns of \( X \) are linearly independent.
The corollary also provides a simple check
for primal uniqueness given a primal-dual solution pair.

\begin{restatable}{lemma}{uniqueCGL}\label{lemma:unique-cgl}
    Fix \( \lambda > 0 \).
    The solution to CGL problem is unique
    if and only if \( \cbr{\Xbi \vi}_{\calS_\lambda} \)
    are linearly independent.
\end{restatable}

Note that a dual solution \( \rho \) is necessary to compute \( v \) in general;
By uniformizing over \( \vi \), we obtain a stronger condition that can be checked
whenever \( \equi \) is known, yet is still
weaker than linear independence of the columns
of \( X \).

\begin{corollary}
    If the columns of \( X_\equi \) are linearly independent,
    then the CGL problem has a unique solution.
\end{corollary}

Finally, we consider the special case when there are no constraints and
\( K = 0 \).
In this setting, \( \vi = \ci \) --- the dual parameters are trivially zero
--- and we can provide a global condition which is much stronger than linear
independence.

\begin{restatable}{proposition}{GGP}[Group General Position]\label{prop:unique-sol}
    Suppose for every \( \calE \subseteq \calB \), \( \abs{\calE} \leq n +1 \),
    there do not exist unit vectors \( \zi \in \R^{\abs{\bi}} \)
    such that for any \( j \in \calE \),
    \[
        X_{b_j} z_{b_j} \in
        \text{affine}(\cbr{\Xbi \zi : \bi \in \calE \setminus b_j}).
    \]
    Then the group lasso solution is unique for all \( \lambda > 0 \).
\end{restatable}
We call this uniqueness condition \emph{group general position}
(GGP) because it naturally extends general position to groups of vectors.
General position is itself an extension of affine independence and is
sufficient for the lasso solution to be unique \citep{tibshirani2013unique}.
GGP is strictly weaker than linear independence of the columns of \( X \),
but neither implies nor is implied by general position (\cref{prop:general-position-relation}).

\subsection{Computing Dual Optimal Parameters}

The main difficulty of \cref{lemma:unique-cgl}
is that knowledge of a dual optimal parameter is required to check if a
unique solution exists.
A dual optimal parameter is also required to fully leverage our
characterization of the optimal set.
As such, now we turn to computing optimal dual parameters.

We give one Lagrange dual problem for CGL in \cref{lemma:lagrange-dual}.
A nice feature of this dual problem is \( \vi \) attains an
alternative interpretation as dual variable.
However, evaluating the dual requires computing \( (X^\top X)^+ \),
which may be difficult even if \( X \) is structured, as in the convex
ReLU program.
Instead, we focus on computing \( \rho \) given a primal solution.

Let \( \w \in \solfn(\lambda) \).
If \( \wi \neq 0 \), then KKT conditions imply
\begin{equation}\label{eq:dual-active}
    \Ki \ri = \ci - \lambda \frac{\wi}{\norm{\wi}_2},
\end{equation}
so that the ``dual fit'' \( \hat d_\bi = \Ki \ri \) is easily computed.
Recovering the dual parameter is a linear feasibility problem:
\begin{equation}\label{eq:dual-active-lp}
    \ri \in \cbr{\ri \geq 0 : \Ki \ri = \hat d_\bi}.
\end{equation}
If \( \wi = 0 \), then complementary slackness is trivially satisfied and we
compute the min-norm dual parameter by solving the following program:
\begin{equation}\label{eq:dual-inactive}
    \min \cbr{\norm{\ri}_2 : \norm{\ci - \Ki \ri}_2 \leq \lambda, \ri \geq 0 }.
\end{equation}
In general, however, we only need \emph{some} dual optimal parameter for
our results to hold; thus, is is typically easier to find \( \rho \)
by solving the following non-negative regression:
\begin{equation}\label{eq:non-negative-regression}
    \ri = \argmin \cbr{\norm{\Ki \ri - \ci}^2_2 : \ri \geq 0 }.
\end{equation}
See \cref{prop:non-negative-regression} for details.

\subsection{Minimal Solutions and Optimal Pruning}\label{sec:pruning}

Often we want the most parsimonious solution, i.e. the one
using the fewest feature groups.
We say a primal solution \( w \) is \emph{minimal} if there
does not exist \( w' \in \solfn(\lambda) \) such that
\( \act(w') \subsetneq \act(w) \).
Building on the previous section, we start with a sufficient
condition for \( w \) to be minimal.
\begin{restatable}{proposition}{minimalSols}\label{prop:minimal-solutions}
    For \( \lambda > 0 \), \( w \in \solfn(\lambda) \) is minimal if and only
    if the vectors \( \cbr{\Xbi \wi}_{\calA(w)} \) are linearly independent.
\end{restatable}

Linear independence of \( \cbr{\Xbi \wi}_{\calA(w)} \)
also identifies \( w \) as a vertex of \( \solfn(\lambda) \)
\citep{bertsekas2009convex},
meaning minimal models are exactly the extreme points of the optimal set.
Combining this characterization with our condition
for uniqueness of a solution (\cref{lemma:unique-cgl}) shows that
minimal solutions are the only solution on their support.

\begin{corollary}\label{cor:minimal-unique}
    Suppose \( w \) is a minimal solution.
    Then \( w \) is the unique solution with support \( \act(w) \).
\end{corollary}

If \( X \) satisfies \emph{group dependence}
(\cref{def:group-dependence}), then all minimal solutions have the same
number of active blocks.

\begin{restatable}{proposition}{minimalSolCharacterization}\label{prop:minimal-sol-characterization}
    Let \( \calV \!=\! \Span(\cbr{\Xbi \bar \wi}) \) for \( \bar w \in \solfn(\lambda) \).
    If \( X \) satisfies group dependence, then every minimal solution has
    \( c = \text{dim}(\calV) \) active blocks.
\end{restatable}

\begin{algorithm}[tb]
	\caption{Optimal Solution Pruning}
	\label{alg:pruning-solutions}
	\begin{algorithmic}
		\STATE {\bfseries Input:} data matrix \( X \), solution \( w \).
		\STATE \( k \gets 0 \).
		\STATE \( w^k \gets w \).
		\WHILE {\( \exists \beta \neq 0 \) s.t. \( \sum_{\bi \in \act(w^k)} \beta_\bi \Xbi \wi^k = 0 \)}
		\STATE \( \bi^k \gets \argmax_{\bi} \cbr{|\beta_\bi| : \bi \in \act(w^k)}  \)
		\STATE \( t^k \gets 1/|\beta_{\bi^k}| \)
		\STATE \( \w^{k+1} \gets \w^k (1 - t^k \beta_\bi) \)
		\STATE \( k \gets k + 1 \)
		\ENDWHILE
		\STATE {\bfseries Output:} final weights \( \w^k \)
	\end{algorithmic}
\end{algorithm}




\cref{alg:pruning-solutions} gives a procedure which, starting from
any optimal solution \( w \), computes an optimal model with minimal
support in \( O((n^3 l + nd) \) time, where \( l \)
is the number active blocks in \( w \)
(see \cref{prop:pruning-correctness}).
Moreover, if the blocks of \( X \) satisfy a regularity condition, then
starting model has no affect on the cardinality of the minimal solution found.
Our algorithm can also be used to verify a minimal solution, since
if \( w \) minimal then it is unique on its support and
\cref{alg:pruning-solutions} must return \( w \) immediately.
This procedure also implies the existence of at least one minimal
solution.

\begin{corollary}\label{cor:smallest-solution-cgl}
    There exists \( w \in \solfn(\lambda) \) for which the vectors
    \(
    \cbr{\Xbi w(\lambda) : \bi \in \calA(w) }
    \)
    are linearly independent.
\end{corollary}

\cref{cor:smallest-solution-cgl} will be useful tool later when we study
sensitivity of the model fit to perturbations in \( y \) and \( \lambda \).

A disadvantage of \cref{alg:pruning-solutions} is that it cannot
continue beyond a minimal solution.
However, minimal models may still be quite large.
We can perform approximate pruning in such cases
using the least squares fit to approximate \( \beta \),
\[ \tilde \beta = \argmin \norm{A \beta - X_{b_j} w_{b_j}}_2^2, \]
where \( A = [\Xbi \wi]_{\act \setminus b_j} \) and \( b_j \in \act \) is
chosen randomly.
Using \( \tilde \beta \) in
\cref{alg:pruning-solutions}
is optimal when \( \cbr{\Xbi \wi}_{\act} \) are dependent
and chooses the update parameters to minimize degradation
of the model fit otherwise.

\subsection{Continuity of the Solution Path}

A major concern when learning with regularizers is how to tune the
parameter \( \lambda \).
Typical strategies like grid-search on the (cross) validation
loss are effective only if the solution function satisfies basic continuity
properties.
For example, if \( \solfn \) is single-valued but discontinuous in \( \lambda \),
then the sample complexity of grid-search can be made arbitrarily poor
by ``hiding'' the optimal \( \lambda \) in a
discontinuity \citep[Sec. 1.1]{nesterov2018lectures}.
In this section, we justify grid-search for CGL by proving several
continuity properties of the solution function, particularly when the
solution is unique.
We start with basic definitions of continuity for point-to-set maps.

\begin{definition}[Closed]
    \( T : \calX \into 2^\calZ \) is closed if
    \( \cbr{\xk} \subset \calX \), \( \xk \into \bar x \) and \( z_k \in T(\xk) \),
    \( z_k \into \bar z \) implies \( \bar z \in T(\bar x) \).
\end{definition}

\begin{definition}[Open]
    \( T : \calX \into 2^\calZ \) is open if
    \( \cbr{\xk} \subset \calX \), \( \xk \into \bar x \) and \( \bar z \in T(\bar x) \),
    implies there exists \( k' \in \bbN \), \( z_k \in T(\xk) \) for \( k \geq k' \),
    such that \( z_k \into \bar z \).
\end{definition}

We say that \( T \) is \emph{continuous} if it is both closed and open.
If \( T(x) \) is a singleton for all \( x \in \calX \),
then openness/closedness are equivalent and imply continuity.
We start with (functional) continuity of the optimal objective.

\begin{restatable}{proposition}{valueContinuityCGL}\label{prop:value-continuity-cgl}
    \( \lambda \mapsto p^*(\lambda) \) is continuous for all
    \( \lambda \geq 0 \).
\end{restatable}

While standard sensitivity results imply that \( \solfn \) is closed,
unfortunately openness is not possible in the general setting.

\begin{restatable}{proposition}{mapContinuityCGL}\label{prop:map-continuity-cgl}
    While \( \solfn \) is closed on \( \R_+ \),
    it is open if only if \( X \) is full column rank.
    However, if the solution is unique on
    \( \Lambda \subset \R^+ \), then \( \solfn \)
    is open at every \( \lambda \in \Lambda \).
\end{restatable}

As a corollary of \cref{prop:map-continuity-cgl}, \( \solfn \) is open
on \( \R^+ \) if and only if \( X \) is full column rank.
Continuity of \( \solfn \) is impossible in general because,
as \citet{hogan1973point} shows, openness is a local stability property;
since \( \solfn(0) \) is unbounded, many ``unstable'' solutions
exist at \( \lambda = 0 \) which are not limit points of other solutions.
Continuity of the unique solution path is an immediate corollary of
\cref{prop:map-continuity-cgl}.

\begin{corollary}
    If the CGL solution is unique on an interval \( \Lambda \subset \R_+\),
    then it is also continuous on \( \Lambda \).
\end{corollary}

In particular, if \( K = 0 \) and GGP holds, then the group lasso solution
is continuous for all \( \lambda > 0 \).
We can strengthen our continuity results when \( \Ki = 0 \) in another way:
by analyzing the dual of the group lasso problem,
we extend continuity from \( p^* \) to the optimal model fit.

\begin{restatable}{proposition}{modelFitContinuity}\label{prop:model-fit-continuity}
    If \( K = 0 \),
    then \( \hat y(\lambda)  \) is continuous on \( \R_+ \) and the penalty
    \( \sum_{\bi \in \calB} \norm{\wi(\lambda)}_2 \) is continuous
    for \( \lambda > 0 \).
\end{restatable}

\subsection{The Min-Norm Path}\label{sec:min-norm}

Now we turn our attention to the min-norm solution path.
Min-norm solutions are typically used in under-determined problems since the
norm of the solution is connected to generalization
\citep{neyshabur2015bias, gunasekar2017implicit}.
Furthermore, the min-norm solution is
a function of \( \lambda \), unlike \( \solfn \).
Throughout this section, \( \acts = \act(\wmin) \) denotes
the active set of the min-norm solution.

Unfortunately, studying the min-norm path immediately encounters a
surprising difficulty:
as opposed to least-squares problems or the lasso (see \citet{tibshirani2013unique}),
the min-norm solution may not lie in the row space of the active set.

\begin{restatable}{proposition}{rowCE}\label{prop:row-ce}
    Suppose \( \Ki = 0 \).
    There exists \( (X, y) \) and
    \( \lambda > 0 \) such that
    \( \wmin_\acts(\lambda) \not \in \Row(\Xas) \).
\end{restatable}

Since the min-norm solution is not given by projecting onto \( \Row(\Xas) \),
how can we compute and study it?
Again, our characterization for the solution set provides a way forward.

\begin{restatable}{proposition}{minNormProgram}\label{prop:min-norm-program}
    Let \( \lambda > 0 \)
    and consider the program:
    \begin{equation}\label{eq:min-norm-program}
        \begin{aligned}
            \alpha^* = \argmin_{\alpha \geq 0} \norm{\alpha}_2^2
            \, \, \text{ s.t.} \, \,
            \sum_{\bi \in \calS_\lambda} \alpha_{\bi} \Xbi \vi = \hat y.
        \end{aligned}
    \end{equation}
    Then the min-norm solution is given by
    \( \wmin_\bi = \alpha^*_\bi \vi \).
\end{restatable}

\cref{eq:min-norm-program} is a quadratic program (QP) that can be solved with
off-the-shelf software like \textsc{cvxpy} \citep{diamond2016cvxpy}.
If \( |\calB| \) and \( d \) are large, this QP may be too expensive to handle directly.
In such situations, we propose to solve the following elastic-net-type problem
\begin{equation}\label{eq:l2-penalized-cgl}
    \begin{aligned}
        \min_{w} \,
         & \half \norm{X w - y}_2^2 + \lambda \sum_{\bi \in \calB} \norm{\wi}_2
        + \frac{\delta}{2} \norm{w}_2^2                                                 \\
         & \quad \text{s.t.} \quad \Ki^\top \wi \leq 0  \text{ for all } \bi \in \calB.
    \end{aligned}
\end{equation}

This \( \ell_2 \)-penalized CGL problem is equivalent to
CGL with modified dataset \( (\tilde X, \tilde y) \) (\cref{lemma:l2-reformulation}).
Since the optimization problem is strongly convex (\( \tilde X \) is full column rank),
invoking \cref{prop:map-continuity-cgl} implies the solution \( w^\delta(\lambda) \) is continuous for all \( \lambda \geq 0 \).
Moreover, as \( \delta \into 0 \), the penalized solution converges
to the min-norm solution to CGL.
\begin{restatable}{proposition}{penConvergence}\label{prop:l2-convergence}
    The solution to the \( \ell_2 \)-penalized problem converges to
    the min-norm solution as \( \delta \rightarrow 0 \).
    That is,
    \[
        \lim_{\delta \into 0} w^\delta(\lambda) = \wmin(\lambda).
    \]
\end{restatable}

Uniqueness and continuity of the solution path for the penalized CGL
problem mean we may prefer to solve
Problem~\eqref{eq:l2-penalized-cgl} with small \( \delta > 0 \) when
tuning \( \lambda \).
\cref{prop:l2-convergence} guarantees that the bias induced by \( \delta \)
will be small and a polishing step with \( \delta = 0 \) can always be used.
Finally, non-zero \( \delta \) ensures the objective is strongly convex,
meaning we can use linearly convergent methods to solve the problem.

%Finally, using stationarity of the Lagrangian gives the following
%closed-form expression for the active blocks of \( w^{\delta}(\lambda) \)
%given \( \vi \):
%\begin{equation}\label{eq:penalized-cfs}
%    w^{\delta}_\acts = \rbr{\Xas^\top \Xas + \delta I}^{-1} \rbr{\Xas^\top y - \vi}.
%\end{equation}

\subsection{Sensitivity}

Now we move onto the problem of sensitivity of a solution
\( w \in \solfn(\lambda, y) \)
to perturbations, either in \( \lambda \) or the targets \( y \).
The main tool for measuring such perturbations are the gradients, for example
\( \nabla_\lambda w(\lambda, y) \).
However, since the solution path of the group lasso is non-smooth,
we must cope with the fact that gradients are not available everywhere.

We show that the gradients of minimal solutions exist
almost everywhere under additional constraint qualifications (CQs).
We do so by considering a \emph{reduced problem} and showing that the
solution to this reduced problem is exactly \( \wa \).
Define the reduced problem as follows:
\begin{equation}\label{eq:constraint-reduced}
    \begin{aligned}
        \min_{\wa} \,
         & \half \norm{\Xa \wa - y}_2^2 + \lambda \sum_{\bi \in \act} \norm{\wi}_2 \\
         & \quad \text{s.t.} \quad \Ka^\top \wa \leq 0
    \end{aligned}
\end{equation}
If \( \act(w) \) is the support of a minimal solution, then \( w \) is the only
solution with support \( \act \) and \cref{eq:constraint-reduced} can be used
to compute the unique active weights.
\begin{restatable}{proposition}{constraintReduced}\label{prop:constraint-reduced}
    Let \( w \in \solfn(\lambda, y) \) be minimal.
    The active blocks \( \wa \) are the unique solution to
    Problem~\eqref{eq:constraint-reduced}.
\end{restatable}
We use this fact to obtain a local solution function
for CGL using the implicit function theorem.
Given a solution \( w \), let
\[
    \calB(w) = \bigcup_{\bi \in \act} \cbr{j \in [a_\bi]: [\Ki]_j^\top \wi = 0 },
\]
be the active constraints.
We now need two classical CQs.

\begin{definition}[LICQ]
    We say that \( w \! \in \! \solfn(\lambda, y) \) satisfies the linear
    independence constraint qualification (LICQ) if
    \( \cbr{[K]_j : j \! \in \! \calB(w)} \) are linearly independent.
\end{definition}

\begin{definition}[SCS]
    Primal solution \( w \in \solfn(\lambda) \) satisfies
    strict complementary slackness if there exists a dual optimal parameter
    \( \rho \) such that \( [\rho]_j > 0 \) for every \( j \in \calB \).
\end{definition}

Now we can state our main differential sensitivity result.

\begin{restatable}{proposition}{localSolFn}\label{prop:local-sol-fn}
    Let \( w \in \solfn(\bar \lambda, \bar y) \) be minimal and suppose
    \( w \) satisfies LICQ on the active set \( \act \)
    and SCS on the equicorrelation set \( \equi \).
    Then \( w \) has a locally continuous solution function
    \( (\lambda, y) \mapsto w(\lambda, y) \).
    Moreover, if
    \[
        D =
        \begin{bmatrix}
            \Xa^\top \Xa + M(\bar w) & \Ka                            \\
            \bar \ra \odot \Ka       & \text{diag}(\Ka^\top \bar \wa)
        \end{bmatrix},
    \]
    where \( \odot \) is the element-wise product,
    \( u_{b_i} = \frac{\wi}{\norm{\wi}_2} \),
    \( u \) is the concatenation of these vectors,
    and \( M \) is block-diagonal projection matrix in \cref{eq:projection-matrix-block},
    then the Jacobians of \( w(\bar \lambda, \bar y) \)
    with respect to \( \lambda \)
    and \( y \) are given as follows:
    \[
        \nabla_\lambda w(\bar \lambda, \bar y)
        = - [D^{-1}]_\act u_\act
        \hspace{0.5em}
        \nabla_y w(\bar \lambda, \bar y)
        = [D^{-1}]_\act \Xa^\top,
    \]
    where \( [D_\act^{-1}]_\act \) is the \( |\act| \times |\act| \)
    dimensional leading principle submatrix of \( D \).
\end{restatable}
